\chapter{Probability, estimation, and valid inference}
A random variable is a measurable function on a probability space
$(\Omega,\mathcal F,\Prob)$. An expectation is an integral with
respect to $\Prob$. The distinction between a deterministic
signal and a stochastic process must be preserved even when both are
stored as the same array of samples.

\entry{40}{Bayes' theorem and total probability}
{For events with positive conditioning probability,
$\Prob(A\mid B)=\Prob(B\mid A)\Prob(A)/\Prob(B)$;
for a partition $(A_i)$, $\Prob(B)=\sum_i\Prob(B\mid A_i)\Prob(A_i)$.}
{Build probabilistic decoders and posterior state estimates.}
\entry{41}{Law of total expectation}
{For integrable $Y$, $\E[Y]=\E[\E[Y\mid G]]$.}
{Marginalize trial-to-trial latent gain or unobserved state.}
\entry{42}{Law of total variance}
{For square-integrable $Y$,
$\Var(Y)=\E[\Var(Y\mid G)]+\Var(\E[Y\mid G])$.}
{Explain Poisson noise plus variable excitability.}
\entry{43}{Law of total covariance}
{For square-integrable $X,Y$,
$\Cov(X,Y)=\E[\Cov(X,Y\mid G)]
+\Cov(\E[X\mid G],\E[Y\mid G])$.}
{Separate residual noise correlation from shared modulating factors.}
\entry{44}{Law of large numbers}
{Sample means converge to the expectation under appropriate
independence or dependence, moment, and ergodicity assumptions.}
{Explain consistency and why many time bins need not imply many independent samples.}
\entry{45}{Central limit theorem}
{For iid observations with finite nonzero variance, normalized
sample-mean fluctuations converge in distribution to a normal law;
dependent-data versions need stronger conditions.}
{Motivate uncertainty quantification for population statistics.}
\entry{46}{Birkhoff ergodic theorem}
{For an integrable observable of a measure-preserving system,
time averages converge almost surely to a conditional invariant
expectation; ergodicity makes the limit equal the ensemble expectation.}
{Clarify when one long recording can replace independent replicates.}
\entry{47}{Markov and Chebyshev inequalities}
{For $X\geq0$, $\Prob(X\geq a)\leq\E[X]/a$;
for finite variance,
$\Prob(|X-\E X|\geq a)\leq\Var(X)/a^2$.}
{Derive conservative uncertainty bounds for estimators.}
\entry{48}{Hoeffding and Bernstein concentration bounds}
{Bound sums of independent bounded variables, or those obeying
specified moment constraints; extensions to dependence require care.}
{Control finite-sample false alarms under valid independence assumptions.}
\entry{49}{Delta method}
{A differentiable map of an asymptotically normal estimator
is asymptotically normal after suitable scaling.}
{Approximate uncertainty of nonlinear indices.}
\entry{50}{Cramer--Rao inequality}
{With regular parametric likelihood and unbiased estimator,
variance is bounded by reciprocal Fisher information
(or its multivariate matrix counterpart).}
{Compare theoretical accuracy limits of firing-rate estimators.}
\entry{51}{Neyman--Pearson lemma}
{For testing one simple hypothesis against another,
thresholding the likelihood ratio yields a most-powerful
test at fixed size, allowing randomization at a threshold.}
{Understand what an optimal detection test does and does not guarantee.}
\entry{52}{Wilks' theorem}
{Under regularity and appropriate interior-point assumptions,
twice a nested-model log-likelihood difference is asymptotically
chi-squared with degrees of freedom equal to parameter difference.}
{Test nested GLM or autoregressive restrictions when assumptions apply.}
\entry{53}{Fisher--Neyman factorization theorem}
{In dominated statistical families, a statistic is sufficient
precisely when the likelihood has an appropriate factorization.}
{Justify reductions of data into counts only for particular models.}
\entry{54}{Exact randomization testing under exchangeability}
{If the null distribution is invariant under a finite transformation
group, a correctly ranked test statistic gives a valid finite-sample
randomization test.}
{Build scientifically matched permutation and jitter controls.}
\entry{55}{Benjamini--Hochberg FDR control}
{For valid $p$-values under independence or certain positive
dependence assumptions, the step-up procedure controls
expected false discovery proportion.}
{Correct across electrodes and spectral comparisons.}
\entry{56}{Isserlis' or Wick's Gaussian moment theorem}
{Centered jointly Gaussian even-order moments decompose into
sums over pairings of covariances; odd-order moments vanish.}
{Evaluate Gaussian baseline moments and covariance estimators.}

\begin{example}[A modulated Poisson count]
Let $N\mid G\sim\Poisson(\mu G)$, $G>0$, and $\E G=1$.
The law of total variance gives
\[
\E N=\mu,\qquad
\Var(N)=\mu+\mu^2\Var(G),\qquad
\frac{\Var(N)}{\E N}=1+\mu\Var(G).
\]
This derivation explains overdispersion without claiming that
a particular dataset actually follows this latent-gain model.
See Goris et al.~\cite{goris2014}.
\end{example}
\begin{caution}
A confidence interval does not automatically have its nominal coverage
after choosing frequency bands, time windows, and hypotheses from
the same data without accounting for selection.
\end{caution}
% --- expanded edition (October 2026) ---
\section{Proofs of selected inference results}
\begin{proof}[Proof of Result 43 (total covariance)]
Write $X-\E X=(X-\E[X\mid G])+(\E[X\mid G]-\E X)$ and similarly for $Y$.
Expanding the product, the two cross terms vanish because
$\E[(X-\E[X\mid G])\,h(G)]=0$ for any square-integrable $h(G)$. The
remaining terms are $\E[\Cov(X,Y\mid G)]$ and
$\Cov(\E[X\mid G],\E[Y\mid G])$.
\end{proof}
For two neurons modulated by a common gain $G$ with
$X\mid G\sim\Poisson(\mu_1G)$, $Y\mid G\sim\Poisson(\mu_2G)$, conditionally
independent, the first term vanishes and $\Cov(X,Y)=\mu_1\mu_2\Var G$.
Shared gain alone produces positive noise correlations proportional to
the product of the rates, which is one reason correlations scale with
firing rate.

\begin{proof}[Proof of Result 50 (Cram\'er--Rao, scalar case)]
Let $U=\partial_\theta\log p_\theta(X)$. Under conditions allowing
differentiation under the integral, $\E_\theta U=0$ and, for any unbiased
$\hat\theta$, $\partial_\theta\E_\theta\hat\theta=1=\E_\theta[\hat\theta U]
=\Cov_\theta(\hat\theta,U)$. Cauchy--Schwarz gives
$1\le\Var_\theta(\hat\theta)\,\Var_\theta(U)=\Var_\theta(\hat\theta)\,\mathcal I(\theta)$.
\end{proof}
For $n$ independent Poisson counts in windows of length $\Delta$ with
rate $r$, $\mathcal I(r)=n\Delta/r$, so no unbiased estimator of the rate
has variance below $r/(n\Delta)$, which the pooled rate attains. More
recording time, not finer binning, is the only way to improve it.

\begin{example}[Result 53: counts are sufficient for Poisson rates]
For a homogeneous Poisson process on $[0,T]$, the likelihood of spike
times $t_1,\ldots,t_n$ is $\lambda^ne^{-\lambda T}$. It depends on the
data only through $n$, so the count is sufficient for $\lambda$: spike
timing carries no further information about the rate \emph{under this
model}. For an inhomogeneous or history-dependent model, the likelihood
depends on the times, and binning to counts discards information.
Sufficiency is always relative to a model.
\end{example}

\section{The delta method in neural statistics}
If $\sqrt n(\hat\theta-\theta)\Rightarrow\Normal(0,\sigma^2)$ and $g$ is
differentiable at $\theta$ with $g'(\theta)\ne0$, then
$\sqrt n(g(\hat\theta)-g(\theta))\Rightarrow\Normal(0,g'(\theta)^2\sigma^2)$.
\begin{proof}[Proof sketch]
By Taylor's theorem $g(\hat\theta)-g(\theta)=g'(\tilde\theta)(\hat\theta-\theta)$
with $\tilde\theta$ between the two; $\tilde\theta\to\theta$ in
probability, so Slutsky's lemma applies.
\end{proof}
\begin{example}[Standard error of power in decibels]
If a spectral estimate is $\hat S\approx S\chi^2_{2K}/(2K)$, then
$\Var\hat S\approx S^2/K$. With $g(S)=10\log_{10}S$,
$g'(S)=10/(S\ln10)$ and
$\operatorname{sd}(10\log_{10}\hat S)\approx 10/(\ln10\sqrt K)\approx4.34/\sqrt K$\,dB,
independent of $S$. The log transform stabilizes the variance, which is
why spectra and their condition differences are analysed in decibels.
\end{example}
\begin{example}[Fano factor]
For $n$ counts with mean $\mu$, variance $v$ and Fano factor $F=v/\mu$,
the multivariate delta method gives an approximate standard error that
grows as the counts become more skewed. A simple bootstrap over trials is
often more reliable than the asymptotic formula when $n<50$.
\end{example}

\section{Concentration: Hoeffding's inequality}
\begin{theorem}[Hoeffding]
If $X_1,\ldots,X_n$ are independent with $a_i\le X_i\le b_i$, then for $t>0$,
\[
\Prob\bigl(\bar X-\E\bar X\ge t\bigr)\le\exp\Bigl(-\frac{2n^2t^2}{\sum_i(b_i-a_i)^2}\Bigr).
\]
\end{theorem}
For a decoder evaluated on $n$ independent test trials with $0/1$ loss,
the observed accuracy is within $t$ of its expectation with probability
at least $1-2e^{-2nt^2}$. With $n=100$ trials and $t=0.1$ this is $0.73$;
to guarantee $\pm0.05$ with probability $0.95$ requires $n\ge738$.
The guarantee assumes independent test trials; neighbouring time bins
from the same trial are not.

\section{Isserlis' theorem and the variance of covariance estimates}
For a centered Gaussian vector,
$\E[X_1X_2X_3X_4]=\Sigma_{12}\Sigma_{34}+\Sigma_{13}\Sigma_{24}+\Sigma_{14}\Sigma_{23}$.
Applied to the product $X_iX_j$, it gives
$\Var(X_iX_j)=\Sigma_{ii}\Sigma_{jj}+\Sigma_{ij}^2$, so the sample
covariance from $T$ independent samples has variance
$(\Sigma_{ii}\Sigma_{jj}+\Sigma_{ij}^2)/T$. Translated into correlations,
the standard error of a near-zero correlation is about $1/\sqrt T$:
detecting pairwise correlations of $0.02$ at two standard errors needs
$T\approx10^4$ independent samples.

\section*{Exercises for this chapter}
\begin{exercise}
Derive the Fisher information of an exponential interspike-interval model
and the Cram\'er--Rao bound for its rate.
\end{exercise}
\begin{exercise}
Use the delta method to find the approximate variance of
$\log(\hat r_1/\hat r_2)$ for two independent Poisson rate estimates.
\end{exercise}
\begin{exercise}
Prove Isserlis' formula for four jointly Gaussian variables by
differentiating the moment generating function $\exp(\tfrac12 s^\top\Sigma s)$.
\end{exercise}
\begin{exercise}
Show by an example that pairwise-independent Bernoulli variables can
violate Hoeffding's bound for their mean.
\end{exercise}
